\section{Discussion}
Reliability-aware scheduling changes where exact calculations are spent without changing what constitutes a physical result.  This separation is essential near shallow thermodynamic minima, momentum-window edges, and gap closings, where a surrogate prediction alone cannot establish a phase.  The two-dimensional calculations show that a small fraction of the prescribed candidate grid can recover the principal thermodynamic structure, and the topology-resolved calculation shows how spectral applicability can be layered onto that result without conflating a gap-closing diagnostic with a stable phase.  The three-dimensional extension demonstrates the value of selective exact sampling as the parameter domain grows, while its discrete support sets a clear limit on the claim.

Exploratory learned-acquisition studies trade local topology support against sequential surprise, but neither tested policy has an overall advantage and no candidate node is validated as a new phase.  These studies motivate future adaptive discovery policies rather than a claim of learned-acquisition superiority.

A longer-term direction is a learned physical world model that represents thermodynamic state, spectra, and calibrated uncertainty jointly.  The present evidence does not yet validate the spectral descriptors or expected-information-gain policy required for such a model, so that direction remains an outlook rather than an established method.
